\documentclass[11pt]{article}
\usepackage[margin=1in]{geometry}
\usepackage{amsmath,amssymb}
\usepackage{times}
\usepackage{hyperref}
\hypersetup{colorlinks=true,linkcolor=blue,urlcolor=blue,citecolor=blue}
\usepackage[numbers,sort&compress]{natbib}
\usepackage{titlesec}
\titleformat{\section}{\normalfont\large\bfseries}{\thesection}{1em}{}
\titleformat{\subsection}{\normalfont\normalsize\bfseries}{\thesubsection}{1em}{}

\title{\textbf{Epistemic Transport and Its Failure: Fragmentation, Migration, and the Loss of Common Epistemic Authority}}
\author{Flyxion \\ Independent Researcher}
\date{September 2026}

\begin{document}
\maketitle

\begin{abstract}
The phrase ``trust apocalypse'' names a real cluster of phenomena, popular disillusionment with science, journalism, medicine, and the courts, alongside the paradoxical growth of demanding traditional institutions amid the decline of looser ones, but the phrase itself is analytically too coarse to explain what it names. It compresses at least three separable processes: a decline in confidence in a given institution across a population, a fragmentation in which different groups converge on incompatible sets of trusted authorities, and a migration of trust away from broad, low-commitment institutions toward narrower, higher-commitment ones. This essay argues that the evidence assembled to support a general collapse of trust more precisely supports a narrower and more tractable claim: that common epistemic authority, the capacity for a claim certified in one social location to remain usable in another, has broken down in identifiable ways, while trust considered locally, within enclaves and dense communities, has not disappeared and may in some cases be intensifying. The essay develops this distinction, calls the capacity that is failing epistemic transport, proposes a recursive-suspicion mechanism that explains why transport specifically fails rather than ordinary skepticism increasing, and closes by applying the resulting framework to a case, public disagreement over vaccine safety, where an individually narratable harm and a statistically reconstructed benefit compete for public credibility under conditions of degraded transport.
\end{abstract}

\section{The Claim That Needs Replacing}

A cluster of observations has recently been gathered under the label ``trust apocalypse,'' developed chiefly by Keiron McCammon in dialogue with John Vervaeke and Jordan Hall, and elaborated by Vervaeke in relation to the co-emergence of religion and cognition: declining confidence in journalism, medicine, and the courts, alongside a puzzling countertrend in which some of the most demanding and doctrinally traditional religious institutions, Eastern Orthodoxy and Traditional Latin Mass Catholicism prominent among them, report renewed interest even as looser, more ecumenical denominations continue to lose adherents \citep{mccammon2025,vervaeke2026}. The label is vivid and has done real work in drawing attention to the pattern. It is nonetheless too coarse to serve as an explanation, for two reasons.

First, ``apocalypse'' suggests a single event, or at least a single underlying process, whereas the phenomena in question are more plausibly the visible effects of several distinct and separable processes operating at different rates in different domains. Second, the word ``trust,'' left unanalyzed, elides a family of related but distinct attitudes, confidence in competence, belief in honesty, willingness to rely on an output despite low affective trust, and perceived legitimacy, that can move independently of one another and that ordinary survey instruments frequently fail to distinguish. A generalized decline in some single quantity called trust is not what the evidence most directly supports. What it supports is examined in the sections that follow.

A further scope limitation should be stated at the outset. Nearly all the quantitative evidence available for this essay concerns the United States. The theoretical claims that follow, about how transport, fragmentation, and recursive suspicion operate, are offered as general hypotheses that a cross-national literature could in principle test, not as findings already established beyond the American case. Where this essay uses phrases like ``a society'' or ``modern societies,'' these should be read as the generalization the American evidence motivates, not as a claim independently verified elsewhere.

\section{Three Processes, Not One}

At least three distinct processes are commonly folded into the single narrative of collapsing trust, and distinguishing them changes what would count as evidence for each.

\emph{Decline} is the simplest process: fewer people across a population express confidence in a given institution than did so at some earlier time. Gallup first measured confidence in a limited set of American institutions in 1973, while its consistently comparable average across nine institutions begins in 1979. That average fell from 48 percent in 1979 to 28 percent in 2025 \citep{gallup2025institutions,gallup2026}. The decline was neither continuous nor uniform. Gallup identifies sharp falls in the early 1980s and early 1990s followed by partial recoveries, and a more durable decline beginning in the mid-2000s \citep{gallup2026}. A separate Gallup series concerning trust in the mass media to report news fully, accurately, and fairly reached 28 percent in 2025, its lowest recorded level \citep{gallup2025media}. These findings establish substantial decline on particular American measures. They do not by themselves establish a simultaneous global collapse of trust across science, journalism, medicine, religion, and law.

\emph{Fragmentation} is a different claim: not that confidence in an institution has fallen for the population as a whole, but that the population no longer converges on which institutions merit confidence, so that trust becomes patterned by prior identity rather than by independent assessment of institutional performance. Kahan and colleagues' work on cultural cognition supplies the relevant mechanism, and Gallup's own party breakdowns supply direct illustration: confidence in the news media among Democrats and Republicans has diverged sharply over the same period in which the population average fell, with the two trends moving toward, and in the most recent data past, opposite ends of the scale rather than converging on a shared low or high level \citep{kahan2011,gallup2025media}. Fragmentation and decline can coexist, but they are not the same finding, and a dataset can show one without showing the other, for instance if two groups both maintain high but mutually exclusive confidence in different sets of institutions, which would show little aggregate decline while showing severe fragmentation.

\emph{Migration} is a third and further claim: that trust has not simply declined or fragmented in place but has relocated, moving away from broad, low-commitment, low-differentiation institutions and toward narrower, higher-commitment, more demanding ones, whether religious, political, or purely informal online communities. The most suggestive case is religious rather than political. Census evidence shows aggregate contraction in major Eastern Orthodox jurisdictions in the United States between 2010 and 2020, a 12 percent membership decline in the Orthodox Church in America and a 22 percent decline in the Greek Orthodox Archdiocese, with substantial variation by diocese and region, while more recent journalistic reporting describes concentrated convert interest in particular parishes, especially among conservative young men attracted to demanding religious practice \citep{krindatch2020,nyt2025orthodox}. These sources do not yet establish the national magnitude or durability of the conversion pattern; a 2025 Pew survey found that only 2 percent of American Catholics attend the Traditional Latin Mass weekly, and older, more optimistic estimates of that movement's size come from self-selected surveys of Latin Mass-affiliated respondents that are not representative of Catholics generally \citep{pew2025catholic}. What the evidence establishes is a more limited conjunction: aggregate decline can coexist with locally intense recruitment into high-commitment communities. This is consistent with migration but does not, by itself, prove it, and the case should be read as suggestive rather than as the clearest available demonstration of the process.

These three processes require different evidence, admit of different causal stories, and call for different remedies. Treating them as a single ``apocalypse'' obscures exactly the distinctions that would make the phenomenon tractable.

\section{What Epistemic Transport Is}

A society does not require its members to hold uniform trust in every institution in order to function epistemically. What it requires is that a claim certified as reliable in one social location remain usable as evidence in another: a laboratory finding that can inform a regulatory decision, a regulatory decision that can inform a news report, a news report that can inform an ordinary conversation, each transfer carrying enough of the original warrant that the claim does not have to be re-derived from scratch by every person who encounters it downstream. Call this capacity epistemic transport.

Transport, not trust in the abstract, is the more precise object of concern. A population can retain substantial local trust, within families, congregations, professional communities, and online affinity groups, while losing the capacity for claims to cross the boundaries between those groups intact. When that capacity fails, a claim certified within one community carries little or no evidential weight in another, not because the second community is simply more skeptical in general, but because it does not recognize the certifying process of the first as itself worthy of deference. The distinctive phenomenon worth naming is not a shortage of trust but a shortage of interoperability between separately maintained pools of it.

\subsection{A Transportable Hypothesis}

The concept needs an operational form or it remains only a suggestive metaphor. Epistemic transport can be operationalized as the change in evidential weight assigned to a claim when its content is held fixed but its certifying source or route of transmission is varied. Let $W_g(C \mid S)$ represent the weight group $g$ assigns to claim $C$ when it is certified by source $S$. A transport failure between groups occurs when a claim accepted on the authority of $S$ within one group undergoes a large credibility discount when presented, with identical content, to another:
\[
T_{g \rightarrow h}(C, S) = W_h(C \mid S) - W_g(C \mid S).
\]
The claim advanced in this essay predicts more than ordinary disagreement about $C$. It predicts that varying only the provenance of an otherwise identical claim will produce large, systematic, and asymmetric changes in acceptance across groups.

This operationalization distinguishes transport failure from generalized skepticism, and the distinction is what keeps the framework falsifiable rather than merely descriptive after the fact. Generalized skepticism predicts uniformly lower confidence across all sources for a given claim. A transport failure predicts source-dependent reversals: the same person may readily accept a claim certified by an in-group source while discounting a substantively identical claim certified elsewhere, and supplying additional evidence through the distrusted channel should have diminishing, null, or even negative effects on acceptance rather than the positive effect ordinary evidential updating would predict. A dataset in which acceptance moves with content regardless of source favors the generalized-decline account; a dataset in which acceptance moves with source almost independently of content favors the transport account. These are, in principle, separable empirical signatures.

\section{Four Attitudes Conflated Under ``Trust''}

Ordinary survey language asking about ``confidence'' in an institution collapses at least four distinguishable attitudes, and the collapse matters because they can move independently.

Competence is a judgment that an institution is technically capable of producing accurate outputs in its domain. Honesty is a judgment that the institution's communications are not deliberately misleading. Reliance is a behavioral disposition to act on an institution's output regardless of affective trust, often because no credible alternative exists; a patient can follow a regulator's dosing guidance while believing the regulator is captured by industry interests, because reliance in the absence of a better option is compatible with low trust in the ordinary sense. Legitimacy is a judgment that an institution's authority to decide a matter is properly constituted, which can survive substantive disagreement with a particular decision, as when a litigant accepts a court's jurisdiction while contesting its ruling.

A single-item confidence measure cannot distinguish a respondent who thinks a regulator is honest but incompetent from one who thinks it competent but dishonest, or a respondent who relies on an institution's output from necessity from one who relies on it from genuine confidence. Claims about a generalized collapse of trust, when traced back to their source, are frequently claims about a decline on one of these four dimensions read as though it applied to all four simultaneously.

\section{Recursive Suspicion and the Failure of Provenance}

The processes described so far explain declining and unevenly distributed confidence. They do not yet explain why transport specifically fails, as opposed to ordinary skepticism simply becoming more common. A further mechanism is needed, and it concerns the object of suspicion rather than its magnitude.

Ordinary skepticism doubts a particular claim and asks for better evidence. A more corrosive pattern doubts the entire chain by which a claim arrived: the researcher's incentives, the institution's funding, the editor's selection criteria, the platform's engagement-optimized distribution, and the listener's own vulnerability to manipulation by all of the above. Because every step in this chain is itself mediated by further institutions, the suspicion recurses rather than terminating in a checkable fact. This is a more specific and more damaging pattern than generalized doubt, because it is not resolved by supplying additional evidence through the same channel; additional evidence delivered by a distrusted process is read as further confirmation that the process is untrustworthy, an availability and amplification dynamic consistent with how risk signals are shown to travel through the ``amplification stations'' that select, repeat, and reinterpret them on the way to a public audience \citep{kasperson1988}. Recursive suspicion of this kind is a plausible mechanism for why claims stop transporting across group boundaries even when the underlying claims have not changed and even when no single step in the chain has become measurably less reliable than before.

This account does not require that institutions have earned no part of the suspicion directed at them. Declining institutional performance, visible error, and a track record of overconfident public communication, examined in the next section, are themselves part of what feeds recursive suspicion; the point is only that suspicion of the whole chain is a distinct phenomenon from suspicion of any single claim, and that it is this whole-chain suspicion, rather than ordinary disagreement about facts, that most directly obstructs transport.

Recursive suspicion is not identical to scrutiny of provenance, and the two should not be conflated. Examining funding sources, incentive structures, editorial selection, and institutional conflicts of interest can be, and often is, a rational condition of warranted belief rather than a pathology. Suspicion becomes transport-destroying specifically when it is both non-terminating and asymmetric: every attempt to verify an out-group claim is disqualified because the verification process itself passes through another distrusted institution, while the equivalent dependencies within a trusted enclave, its own funding, its own selection effects, its own incentives, cease to be noticed or interrogated at all. The pathology, on this account, is not excessive attention to provenance as such. It is the disappearance of any point, on either side of a boundary, at which provenance can be mutually checked, so that scrutiny applies in only one direction across the boundary that recursive suspicion has created.

\section{Causal Pluralism}

An adequate account resists collapsing the causes of declining and fragmenting confidence into a single explanation, because the candidate mechanisms operate at different levels and are not synonyms for one another.

Institutional failure changes the evidence available for judging institutional trustworthiness directly: publicized errors, reversed guidance, and revealed conflicts of interest are genuine updates on institutional performance, not artifacts of perception. Increased visibility changes access to evidence about institutions without necessarily changing institutional performance itself; a failure rate that has not risen can nonetheless generate more publicized failures if reporting and investigative capacity has expanded, or if a platform's incentive structure rewards surfacing failures over surfacing routine competence. Polarization changes how identical institutional conduct is interpreted by different audiences, the mechanism documented in the cultural-cognition literature, so that the same regulatory decision can be read as prudent caution by one group and as proof of capture by another \citep{kahan2011}. Algorithmic curation changes which failures are encountered and how often, independent of their true frequency, by optimizing for engagement rather than representativeness. Declining social capital, in Putnam's sense of the associational density through which institutional claims were historically validated by people one knew personally rather than by anonymous institutions directly, weakens an intermediate layer that once absorbed some of this work \citep{putnam2000}.

These mechanisms interact rather than competing for credit as the single explanation. A given episode of collapsing regional trust in, say, a public-health institution is likely to show contributions from several of them at once, and an essay that reduces the phenomenon to any one, institutional failure alone, or social media alone, is likely to be wrong about the relative weights even where it is right that the mechanism it names is operative.

\section{The Baseline Problem}

Confidence in most major American institutions has been declining, by Gallup's own measure, more or less continuously since the trend began in 1979 \citep{gallup2026}. This creates an immediate problem for any claim that something new and apocalyptic is occurring now: a contemporary ``trust apocalypse'' has to specify what has changed recently that was not already true across four and a half decades of decline, on pain of the label being a dramatic redescription of an old and gradual trend rather than a diagnosis of a new one.

Several candidates are available and are not mutually exclusive. The collapse of information scarcity means that the mere fact of institutional endorsement, once a scarce and therefore informative signal, is now one weak signal among an abundance of freely available competing claims, reducing its marginal evidential value regardless of whether institutions have become less reliable. Algorithmic personalization means that different audiences increasingly encounter different evidentiary environments, weakening the shared factual starting point that earlier eras of mass broadcast media provided even amid disagreement about how to interpret it. The proliferation of alternative, non-institutional sources of apparent expertise, independent podcasters, influencers, and online communities capable of building substantial audiences without passing through any institutional gatekeeper, provides destinations for the migration process described above that did not exist in comparable form even two decades ago. Each of these is a plausible candidate for what is genuinely new; none of them requires positing an acceleration or qualitative break in the underlying decline curve itself, since the decline curve was already present well before any of them existed.

\section{The Religious-Revival Paradox as a Predicted Consequence}

The pattern with which this essay opened, aggregate decline in institutional religious affiliation alongside genuine, if modest, growth in the most doctrinally demanding and traditionally structured religious communities, is often presented as a paradox against a backdrop of a general secularization or trust-collapse narrative. Under the three-process framework developed above, it is not a paradox but a predicted instance of migration.

Henrich's account of credibility-enhancing displays offers a plausible mechanism: cultural learners are not indiscriminate imitators but weight a model's expressed beliefs by the costliness of that model's observable commitment to them, and the CRED framework predicts that institutions capable of eliciting more visibly costly displays of commitment from their members can transmit belief more credibly than those that cannot \citep{henrich2009}. Under conditions of fragmenting trust, demanding institutions may possess an advantage on this account: their practices, fasting, long liturgies, distinctive dress, public doctrinal commitment, make commitment publicly legible through costly and repeated participation in a way that a low-commitment, ecumenical institution structurally cannot easily replicate. Whether this mechanism actually explains the specific pattern of recent Orthodox conversion, however, rather than political sorting, demographic concentration among a particular age and sex cohort, dissatisfaction with the theology or governance of other churches, or selective journalistic attention to an eye-catching subculture, remains an open empirical question that the sources cited here do not settle. Convert accounts of Orthodoxy's appeal do consistently emphasize the tradition's ``rigid, unbending'' structure as an attraction in contrast to institutions perceived as offering only low-cost affiliation, which is at least consistent with the CRED account \citep{nyt2025orthodox}. The paradox with which this section opened is, at minimum, not evidence against the three-process framework: aggregate decline and concentrated high-commitment recruitment are compatible outcomes under migration, whatever the ultimately correct mechanism turns out to be.

\section{Application: Vaccine Testimony as a Transport Failure}

The framework developed here has a direct application to public disagreement over vaccine safety, though it should be read as an application and not as independent evidence for the framework itself.

A person's report of symptoms experienced after vaccination is, in the relevant sense, locally transportable as testimony: the audience can accept that the person suffered without first accepting any regulatory or epidemiological institution's assessment, since the account requires only that the audience believe the speaker experienced what they describe. The further conclusion that vaccination caused those symptoms is not contained in the experience itself; that conclusion requires a causal comparison of exactly the kind Sections 2 and 4 of this essay's companion analysis on vaccine injury statistics develop, though the narrative form of the testimony, arriving as a single continuous story of exposure followed by symptom, can make the causal conclusion feel already established when it is not. A population-level estimate that a given vaccine prevented some very large number of deaths is not locally transportable in the same sense; the claim depends on a transmission model fit to reported and excess mortality across many countries and cannot be independently verified by an individual audience member, so that its persuasive force depends entirely on whether the institutions that produced it, and the further institutions that reported on it, are treated as trustworthy carriers \citep{watson2022}. Where common epistemic authority is intact, this dependence is not a liability, since the modeled estimate simply travels on the credit of the institutions that vouch for it. Where epistemic transport has broken down along the lines described above, the same dependence becomes a liability for the claim's ability to compete with the individually narratable one, not because the modeled estimate is less true but because it has no path to an audience that does not already extend credit to the chain of institutions carrying it.

The differentiated causal evidence this essay's companion piece on vaccine injury statistics develops at length, the distinction between an adverse event and an adjudicated net harm, the specific, real, causally established harms such as the mRNA-myocarditis association, and the non-identifiability of the joint distribution connecting individual harm to individual benefit, is exactly the kind of claim most vulnerable to a transport failure, because every step in its construction, laboratory data, regulatory review, statistical modeling, science journalism, passes through an institutional chain that recursive suspicion can target at any link \citep{nasem2024}. Nocebo research supplies a possible reinforcing pathway rather than the whole story. Placebo recipients in COVID-19 vaccine trials reported substantial rates of systemic adverse events, showing that reported post-vaccination symptoms cannot be attributed wholly to a vaccine's pharmacological action \citep{haas2022}. Prospective and retrospective studies have also found associations among exposure to interpersonal or social-media accounts of side effects, pre-vaccination expectations, and subsequent side-effect experience, with expectations statistically mediating the relationship between social exposure and reported symptoms \citep{tan2022,clemens2023}. These findings are consistent with socially transmitted expectations shaping symptom experience, and with locally transportable testimony outcompeting institutionally mediated analysis for uptake independent of relative accuracy, but they do not establish that online narratives caused any specified increase in clinically validated vaccine injuries, and the causal language should remain restrained accordingly.

\section{Restoring Transport, Not Restoring Trust}

If the diagnosis offered here is broadly correct, the appropriate response is not a demand that the public extend more trust to institutions currently distrusted, since demanding trust does nothing to address whatever has made recursive suspicion of the relevant chain reasonable in the first place, and simply asserting that suspicion is unwarranted does not restore what it targets. O'Neill's distinction between trust and trustworthiness identifies the more tractable target: an institution's task is to become assessably worthy of reliance, through intelligible evidence, accountability for error, and communicative practices that allow a non-expert to evaluate the process rather than simply accept its conclusion, rather than to seek greater trust as an end pursued independently of the practices that would warrant it \citep{oneill2002}. Restoring transport, in the terms developed here, means building points in the institutional chain where an outside party can check a link without having to accept the entire chain on faith, denominators reported alongside numerators, causal standards stated alongside raw counts, corrections issued visibly rather than quietly, and the specific evidentiary basis of a claim kept separable from the institutional reputation of whoever is reporting it. None of this guarantees that transport will be restored; recursive suspicion, once established, does not necessarily yield to better evidence delivered through the same distrusted channels. It does, however, identify the correct object of repair. The phenomenon in question is not that people have stopped trusting. It is that the claims requiring the most trust to transport are, for identifiable and only partly remediable reasons, the ones now least able to complete the journey.

\bibliographystyle{plainnat}
\begin{thebibliography}{99}

\bibitem[Clemens et al.(2023)]{clemens2023}
Clemens, K.~S., Faasse, K., Tan, W., Colagiuri, B., Colloca, L., Webster, R., Vase, L., Jason, E., \& Geers, A.~L. (2023). Social communication pathways to COVID-19 vaccine side-effect expectations and experience. \textit{Journal of Psychosomatic Research}, 164, 111081.

\bibitem[Gallup(2025a)]{gallup2025institutions}
Gallup. (2025a). \textit{Democrats' Confidence in U.S. Institutions Sinks to New Low}. Washington, DC: Gallup, Inc. \url{https://news.gallup.com/poll/692633/democrats-confidence-institutions-sinks-new-low.aspx}

\bibitem[Gallup(2025b)]{gallup2025media}
Gallup. (2025b). \textit{Trust in Media at New Low of 28\% in U.S.} Washington, DC: Gallup, Inc. \url{https://news.gallup.com/poll/695762/trust-media-new-low.aspx}

\bibitem[Gallup(2026)]{gallup2026}
Gallup. (2026). \textit{Confidence in U.S. Institutions Remains Near All-Time Low}. Washington, DC: Gallup, Inc. \url{https://news.gallup.com/poll/712436/confidence-institutions-remains-near-time-low.aspx}

\bibitem[Haas et al.(2022)]{haas2022}
Haas, J.~W., Bender, F.~L., Ballou, S., Kelley, J.~M., Wilhelm, M., Miller, F.~G., Rief, W., \& Kaptchuk, T.~J. (2022). Frequency of adverse events in the placebo arms of COVID-19 vaccine trials: a systematic review and meta-analysis. \textit{JAMA Network Open}, 5(1), e2143955.

\bibitem[Hardwig(1985)]{hardwig1985}
Hardwig, J. (1985). Epistemic dependence. \textit{The Journal of Philosophy}, 82(7), 335--349.

\bibitem[Henrich(2009)]{henrich2009}
Henrich, J. (2009). The evolution of costly displays, cooperation and religion: credibility enhancing displays and their implications for cultural evolution. \textit{Evolution and Human Behavior}, 30(4), 244--260.

\bibitem[Kahan et al.(2011)]{kahan2011}
Kahan, D.~M., Jenkins-Smith, H., \& Braman, D. (2011). Cultural cognition of scientific consensus. \textit{Journal of Risk Research}, 14(2), 147--174.

\bibitem[Kasperson et al.(1988)]{kasperson1988}
Kasperson, R.~E., Renn, O., Slovic, P., Brown, H.~S., Emel, J., Goble, R., Kasperson, J.~X., \& Ratick, S. (1988). The social amplification of risk: a conceptual framework. \textit{Risk Analysis}, 8(2), 177--187.

\bibitem[Krindatch(2020)]{krindatch2020}
Krindatch, A. (2020). \textit{US Religion Census 2020: A Decade of Dramatic Changes in American Orthodox Churches}. Second Census of Orthodox Christian Churches. \url{https://orthodoxreality.org/latest/second-census-2020-of-american-orthodox/}

\bibitem[McCammon et al.(2025)]{mccammon2025}
McCammon, K., Hall, J., \& Vervaeke, J. (2025). Navigating the trust apocalypse: examining collective agency and distributed cognition. \textit{The Lectern}, August 13, 2025.

\bibitem[National Academies of Sciences, Engineering, and Medicine(2024)]{nasem2024}
National Academies of Sciences, Engineering, and Medicine. (2024). \textit{Evidence Review of the Adverse Effects of COVID-19 Vaccination and Intramuscular Vaccine Administration}. Washington, DC: The National Academies Press. \url{https://doi.org/10.17226/27746}

\bibitem[New York Times(2025)]{nyt2025orthodox}
New York Times. (2025). \textit{Orthodox Church Pews Are Overflowing With Converts}. New York: The New York Times Company.

\bibitem[O'Neill(2002)]{oneill2002}
O'Neill, O. (2002). \textit{A Question of Trust: The BBC Reith Lectures 2002}. Cambridge: Cambridge University Press.

\bibitem[Pew Research Center(2025)]{pew2025catholic}
Pew Research Center. (2025). \textit{Catholic Practices and Devotions}. Washington, DC: Pew Research Center. \url{https://www.pewresearch.org/religion/2025/06/16/catholic-practices-and-devotions/}

\bibitem[Putnam(2000)]{putnam2000}
Putnam, R.~D. (2000). \textit{Bowling Alone: The Collapse and Revival of American Community}. New York: Simon \& Schuster.

\bibitem[Tan et al.(2022)]{tan2022}
Tan, W., Colagiuri, B., \& Barnes, K. (2022). Factors moderating the link between personal recounts of COVID-19 vaccine side effects viewed on social media and viewer postvaccination experience. \textit{Vaccines}, 10(10), 1611.

\bibitem[Vervaeke(2026)]{vervaeke2026}
Vervaeke, J. (2026). Cognition, religion, and the trust apocalypse are more connected than you think. \textit{John Vervaeke's Lectern}, September 11, 2026.

\bibitem[Watson et al.(2022)]{watson2022}
Watson, O.~J., Barnsley, G., Toor, J., Hogan, A.~B., Winskill, P., \& Ghani, A.~C. (2022). Global impact of the first year of COVID-19 vaccination: a mathematical modelling study. \textit{The Lancet Infectious Diseases}, 22(9), 1293--1302.

\end{thebibliography}

\end{document}

